% `tex.web` §560 and §561: a `.tfm` that cannot be read is not a font.
%
%     procedure start_font_error_message;
%     begin print_err("Font "); sprint_cs(u); print_char("=");
%     print_file_name(nom,aire,"");
%     if s>=0 then begin print(" at "); print_scaled(s); print("pt"); end
%     else if s<>-1000 then begin print(" scaled "); print_int(-s); end;
%     end;
%     …
%     if file_opened then print(" not loadable: Bad metric (TFM) file")
%     else print(" not loadable: Metric (TFM) file not found");
%
% and `read_font_info` returns `null_font`, so the identifier the document
% asked for means \nullfont from then on, the file name is never registered,
% and a second \font naming it fails again and reports again.
%
% We gave it a slot measuring zero and said nothing, on the theory that one
% missing `.tfm` should cost one wrong box rather than the whole document. But
% it costs the document silently: every glyph in that font vanishes from the
% page with no record anywhere that a font was missing, and \fontname answers
% with a file that was never read.
%
% The other half is §379, reached from §367's `input` case:
%
%     input: if chr_code>0 then force_eof:=true
%       else if name_in_progress then insert_relax
%       else start_input;
%
% A file name is being scanned, so `\input` must not open anything: §379 puts
% it back and a `\relax` in front of it. §526 raises `name_in_progress` over
% the name and §1258 raises it again over the `at`/`scaled` keyword, so it
% catches an `\input` on either side of the size. We ended the name at the
% control sequence and backed it up unexpanded, which reads the same and is
% not — the error context showed `\input` where TeX shows the inserted
% `\relax` above an elided level, and one command went missing from the trace.
%
% Not covered here: "Bad metric (TFM) file", the branch where the file opens
% and its contents are wrong. It needs a malformed `.tfm` on the search path,
% which a `.tex` case cannot put there.
\catcode`\{=1 \catcode`\}=2
\errorcontextlines=5
\showthe\escapechar

% Plain: the name, the message, and \nullfont after it.
\font\gone=mumblezz
\message{PROBE:gone=[\fontname\gone]}

% A second \font naming the same file is not a reuse of the first — nothing
% was registered — so it fails again.
\font\gone2=mumblezz
\message{PROBE:again=[\fontname\gone2]}

% §560 echoes the size spec as it was written.
\font\atsize=mumblezz at 12pt
\font\scal=mumblezz scaled 500
\font\scal1000=mumblezz scaled 1000

% ...and an `\input` in the way, on both sides of the size.
\immediate\openout0=fsub.tex
% The sub-file ends with a report rather than a message, so that TeX's
% closing `)` — written against the current column, and this engine writes no
% file echo at all — lands on a line of its own.
\immediate\write0{\noexpand\message{PROBE:sub}\noexpand\showthe\noexpand\escapechar}
\immediate\closeout0
\font\iname=mumblezz\input fsub
\font\ikey=mumblezz \input fsub

% A font that *is* there still loads, and is still shared with itself.
\font\ok=cmr10
\font\ok2=cmr10
\message{PROBE:ok=[\fontname\ok][\ifx\ok\ok2 same\else different\fi]}
\showthe\escapechar
\end
